\section{Capítulo~\ref{sec:code}. Código Morse}

Las estimaciones límite del capítulo son teoría de la información y cálculo; está medido dónde surge el ruido en el canal, qué tan rápido trabaja el cerebro y cuánto cuesta una espiga, mientras que la pregunta de qué código usa realmente la corteza sigue abierta.

{\footnotesize
\setlength{\tabcolsep}{3pt}\renewcommand{\arraystretch}{1.15}
\begin{longtable}{>{\raggedright}p{0.5cm}>{\raggedright}p{4.0cm}>{\raggedright}p{5.6cm}>{\raggedright}p{2.3cm}>{\raggedright\arraybackslash}p{1.7cm}}
\toprule
N.º & Proposición & Experimento y qué se midió & Fuente & Estatus \\
\midrule\endfirsthead
\toprule
N.º & Proposición & Experimento y qué se midió & Fuente & Estatus \\
\midrule\endhead
1 & Las frecuencias de las neuronas \nt{GLUT} de la corteza van de fracciones a decenas de espigas por segundo, y la mayor parte del tiempo las neuronas trabajan cerca del límite inferior & Registro con una pipeta adosada a la célula (la elección de la neurona no depende de su actividad) en la corteza auditiva de la rata despierta: en cada momento, menos del $5\%$ de las neuronas supera las $20$ espigas por segundo; en parte de las neuronas la frecuencia de fondo es menor de $0.01$ por segundo; la distribución de frecuencias es lognormal & \cite{Hromadka2008} & medido \\
2 & La capacidad del canal de espigas es $R_{\max}\approx r\log_2\frac{e}{r\Delta t}$~\eqref{eq:spikeentropy}, y casi toda está en el tiempo de las espigas (proposición~\ref{prop:capacity}) & Entropía de una cadena binaria de celdas de longitud $\Delta t$. Cifras del capítulo: $8$ bits por espiga con $r=10$ por segundo y $\Delta t=1\ms$, unos $5$ bits con $\Delta t=10\ms$, menos de $2$ bits por segundo para un contador de espigas & \cite{MacKay1952} & teoría \\
3 & Las neuronas sensoriales transmiten de uno a varios bits por espiga, en los mejores casos hasta la mitad del límite & Reconstrucción del estímulo a partir de las espigas de neuronas sensoriales cuyo estímulo se conoce; resumen de mediciones en el libro & \cite{Rieke1997} & medido \\
4 & El generador de espigas es preciso; el tiempo exacto lo fijan los cambios rápidos de la entrada & Neuronas de corteza de rata en rebanadas: con una corriente repetida de fluctuaciones rápidas, las espigas se repiten con precisión mejor que $1\ms$; con corriente constante, los instantes se dispersan & \cite{Mainen1995} & medido \\
5 & La sinapsis es poco fiable: más de la mitad de las espigas no llega al receptor por el azar de la liberación & Estimulación mínima de las entradas de neuronas piramidales del hipocampo (rebanadas): más de la mitad de las espigas no provoca respuesta. Se descartaron las fluctuaciones del umbral de estimulación, los fallos de conducción en las ramificaciones del axón y la baja temperatura del experimento; queda la liberación probabilística & \cite{AllenStevens1994} & medido \\
6 & El flujo de espigas en la corteza viva parece aleatorio: los intervalos se dispersan como en un flujo de Poisson y la varianza del número de espigas es cercana a la media; parte del “ruido” es un mensaje sobre los movimientos del animal & Registros en las áreas visuales V1 y MT del mono despierto: coeficiente de variación de los intervalos cercano a $1$, razón varianza/media del número de espigas como en un proceso aleatorio; un modelo que suma entradas pequeñas independientes da una dispersión mucho menor. En la corteza visual del ratón, una parte considerable de la dispersión se predice a partir de un video de los movimientos del propio animal & \cite{Softky1993,Stringer2019} & medido \\
7 & El código de frecuencia es lento: error $1/\sqrt{rT}$~\eqref{eq:rateerror}, mientras que el cerebro reconoce una imagen en $\sim150\ms$ & La fórmula es estadística de Poisson, deducción del capítulo. Experimento: personas decidían si había un animal en una fotografía desconocida mostrada durante $20\ms$; los potenciales cerebrales para “sí” y “no” divergen al cabo de unos $150\ms$. La división de ese tiempo en una decena de etapas de $10$--$20\ms$ es una estimación del capítulo, no una medición & \cite{Thorpe1996} & medido \\
8 & El promediado sobre un grupo está limitado por la correlación: el error cae como mucho $1/\sqrt c$ veces~\eqref{eq:correlated} (proposición~\ref{prop:pooling}) & Pares de neuronas del área MT del mono registrados a la vez: coeficiente de correlación de los números de espigas de $0.12$ en promedio, y el análisis teórico muestra que incluso esa correlación limita la ganancia del grupo. Teoría: el límite lo pone solo la parte de la correlación que se parece a la señal. Modelo de la postura contraria: la frecuencia de un grupo de $50$--$100$ neuronas se lee en un solo intervalo entre espigas, $10$--$50\ms$, y el tiempo de las espigas individuales no se usa en la corteza & \cite{Zohary1994,MorenoBote2014,ShadlenNewsome1998} & medido \\
9 & Un número puede escribirse en el tiempo de la primera espiga~\eqref{eq:latency}, en el orden de las espigas de un grupo o en la fase de un ritmo local & Retina de salamandra: la estructura espacial de una imagen mostrada brevemente está escrita en las latencias relativas de las primeras espigas de las neuronas de salida, casi independientemente del contraste. Células de lugar del hipocampo de rata: al atravesar “su” lugar, las espigas se adelantan cada vez más en el ciclo del ritmo theta ($7$--$12$\,Hz), con un desplazamiento de $100^\circ$ a $355^\circ$. Hasta qué punto usa la corteza el tiempo de las espigas es una pregunta abierta (fila~8) & \cite{Gollisch2008,OKeefe1993} & medido \\
10 & Qué código se lee lo decide la constante de tiempo del receptor~\eqref{eq:reader} (proposición~\ref{prop:reader}); en la corteza activa, las neuronas están más cerca de los detectores de coincidencias & La fórmula~\eqref{eq:reader} es una deducción del capítulo. Revisión de registros in vivo: en la corteza activa la conductancia de la membrana es varias veces mayor que en reposo, y la constante de tiempo, correspondientemente más corta. Que la corteza cambie de código precisamente así no se ha mostrado directamente & \cite{Destexhe2003} & hipótesis \\
11 & La sinapsis con depresión transmite el cambio de frecuencia, en esencia $\Delta r/r$~\eqref{eq:depression} (fig.~\ref{fig:depression}) & Registros por pares de neuronas piramidales de la corteza: la velocidad de la depresión la fija la probabilidad de liberación; con depresión lenta la respuesta refleja la frecuencia, y con depresión rápida, las coincidencias. Modelo basado en la depresión medida en la corteza de rata: cambios relativos iguales de frecuencia en entradas rápidas y lentas dan la misma respuesta, es decir, la sinapsis transmite $\Delta r/r$. Las cifras $1.7\to2.2$, $6.7$ y $90\ms$ son un cálculo del capítulo & \cite{Tsodyks1997,Abbott1997depression} & medido \\
12 & La ráfaga es una “raya”: su probabilidad de perderse, $(1-p)^k$~\eqref{eq:burst}, es pequeña, y la facilitación la reduce aún más & Revisión: en sinapsis poco fiables, las espigas aisladas se pierden a menudo, mientras que las ráfagas se transmiten de forma fiable gracias a la facilitación; una ráfaga provoca cambios sinápticos duraderos. Que el cerebro lea la ráfaga como un signo aparte es una hipótesis & \cite{Lisman1997,AllenStevens1994} & hipótesis \\
13 & Las neuronas \nt{GABA} rápidas fijan el ritmo y la “puntuación”; otra clase inhibe a las inhibidoras y abre la compuerta (proposición~\ref{prop:punctuation}) & Revisión de las propiedades de las clases de neuronas \nt{GABA} de la corteza. Optogenética: la excitación rítmica de las neuronas \nt{GABA} rápidas provoca el ritmo gamma, y la respuesta a un estímulo depende de la fase del ciclo. En la corteza visual del ratón, las neuronas PV se inhiben entre sí; las SST inhiben a todas las demás clases; las VIP, sobre todo a las SST. La excitación de las neuronas VIP en el ratón despierto retira la inhibición de las neuronas \nt{GLUT}; las activa una señal de refuerzo. La división de papeles como regla general es una hipótesis & \cite{Tremblay2016,Cardin2009,Pfeffer2013,Pi2013} & hipótesis \\
14 & La energía limita la frecuencia media a un valor del orden de una espiga por segundo~\eqref{eq:energy} & Presupuesto energético de la sustancia gris de roedor a partir de datos anatómicos y fisiológicos: las espigas y las corrientes sinápticas suponen el $47\%$ y el $34\%$ de la energía de señalización; no pueden estar activas a la vez más de $\sim15\%$ de las neuronas. Para la corteza humana: no más de $\sim1\%$ de las neuronas pueden estar muy activas a la vez. La cifra de $\sim10^9$ ATP por espiga y la potencia de $\sim10$\,W para la señalización son una estimación del capítulo basada en estos cálculos & \cite{Attwell2001,Lennie2003} & teoría \\
15 & El código es disperso y está ajustado al máximo de bits por julio; la corteza cambia precisión por energía (proposición~\ref{prop:economy}) & Hipótesis del código económico. Apoyo: en ratones con restricción de alimento, en la corteza visual caen la conductancia de los receptores AMPA y el gasto sináptico de ATP ($-29\%$), la frecuencia de espigas se mantiene y la sintonía a la orientación se ensancha un $32\%$ & \cite{Barlow1961,Padamsey2022} & hipótesis \\
\bottomrule
\end{longtable}}
